% !TeX root = ../main.tex
\section{Математическая модель оптимального обнаружителя}
В условиях априорно неизвестной случайной начальной фазы $\varphi_0$ оптимальный обнаружитель строится по квадратурной схеме с некогерентным накоплением. Синфазная $I(\tau_{\text{з}}, f_{\text{д}})$ и квадратурная $Q(\tau_{\text{з}}, f_{\text{д}})$ составляющие коррелятора вычисляются как:
\begin{align}
I(\tau_{\text{з}}, f_{\text{д}}) &= \sum_{k=0}^{N-1} y_k \cdot G(k T_d - \tau_{\text{з}}) \cos\bigl(2\pi(f_0 + f_{\text{д}}) k T_d\bigr), \\
Q(\tau_{\text{з}}, f_{\text{д}}) &= -\sum_{k=0}^{N-1} y_k \cdot G(k T_d - \tau_{\text{з}}) \sin\bigl(2\pi(f_0 + f_{\text{д}}) k T_d\bigr).
\end{align}
Решающая статистика формируется как квадрат огибающей корреляционного отклика:
\begin{equation}
X^2(\tau_{\text{з}}, f_{\text{д}}) = I^2(\tau_{\text{з}}, f_{\text{д}}) + Q^2(\tau_{\text{з}}, f_{\text{д}}).
\end{equation}

\subsection{Вычисление порога обнаружения по критерию Неймана--Пирсона}
При гипотезе $H_0$ (сигнал в проверяемой ячейке отсутствует) синфазная и квадратурная компоненты определяются только шумовой составляющей. Поскольку $n_k$ --- центрированный белый шум с дисперсией $\sigma_n^2 = 1$, величины $I$ и $Q$ представляют собой сумму большого числа независимых случайных величин ($N = 33000$) и по центральной предельной теореме распределены асимптотически нормально с нулевым математическим ожиданием и дисперсией:
\begin{equation}
\begin{split}
\sigma_0^2 = \mathbb{D}[I] = \mathbb{D}[Q] &= \sum_{k=0}^{N-1} \sigma_n^2 \cdot G^2(k T_d - \tau_{\text{з}}) \cos^2\bigl(2\pi(f_0 + f_{\text{д}}) k T_d\bigr) \\
&\approx \frac{N \sigma_n^2}{2} = \frac{33000}{2} = 16500.
\end{split}
\end{equation}
Величина $\frac{X^2}{\sigma_0^2} = \left(\frac{I}{\sigma_0}\right)^2 + \left(\frac{Q}{\sigma_0}\right)^2$ имеет распределение хи-квадрат с двумя степенями свободы ($\chi_2^2$), что эквивалентно экспоненциальному закону распределения для $X^2$:
\begin{equation}
p(X^2 \mid H_0) = \frac{1}{2\sigma_0^2} \exp\left(-\frac{X^2}{2\sigma_0^2}\right), \quad X^2 \ge 0.
\end{equation}
Вероятность ложной тревоги для одной ячейки поиска определяется интегралом:
\begin{equation}
P_{\text{ло}} = \int_{h}^{\infty} p(X^2 \mid H_0) \, d(X^2) = \exp\left(-\frac{h}{2\sigma_0^2}\right).
\end{equation}
Отсюда аналитическое выражение для порога обнаружения $h$:
\begin{equation}
h = -2\sigma_0^2 \ln(P_{\text{ло}}).
\end{equation}
Задавая типовую для задачи первичного поиска вероятность ложной тревоги $P_{\text{ло}} = 10^{-4}$ и подставляя значение дисперсии $2\sigma_0^2 \approx 33000$:
\begin{equation}
h = -33000 \cdot \ln(10^{-4}) \approx 33000 \cdot 9.2103 \approx 3.04 \cdot 10^5.
\end{equation}
С учетом квантования 3-уровневого АЦП эффективная мощность шума несколько снижается, что дает экспериментальную дисперсию шума в ячейках вне пика порядка $1.0\cdot 10^4 \dots 1.5\cdot 10^4$, а соответствующий порог составляет порядка $(1.2 \dots 1.5) \cdot 10^5$.
